\section{A proposed roadmap for groups with device and fabrication resources}\label{sec:roadmap}
The simulations suggest what to measure first. None of the criteria below is an established device qualification: each is
conditional on this simulator (ideal reads, the pulse surrogate, small models) and on the chopped Tiki-Taka setting that passed
in Section~\ref{sec:main}.

\paragraph{Stage 1: single-device characterisation} (tens of devices, no array).
\begin{itemize}
\item \textbf{Measurement.} At least 1000 up and 1000 down pulses at the intended width and voltage, repeated at least three
      times on at least 32 devices, with the read current measured (not inferred) and incomplete pulses rejected. Fit: step
      against state (saturation $\gamma_\pm$), up/down asymmetry, cycle-to-cycle noise (separated from read noise),
      device-to-device spreads, the symmetry point and the signed weight range around it, and drift.
\item \textbf{What the simulations support} (c-TTv2 sheet, CIFAR $c=32$, absolute 2\,pp line, three seeds,
      Table~\ref{tab:rev_spec}): a step asymmetry up to \RVSpecCttvTwoCifarAsymThr{} and a reference error up to
      \RVSpecCttvTwoCifarRefErrThr{} of the half range, one axis at a time around the protonic-inspired base. Thresholds for
      levels, saturation, write noise and retention are not established by this grid; what it shows is where training clearly
      fails (no saturation or $\gamma\ge1$, reference error 0.2, retention of $10^3$ updates or less). A device group should
      therefore measure all of these and simulate its own fitted device, rather than read pass lines off this sheet. Combined
      deviations are not covered.
\item \textbf{Endurance and pulse budget.} Report per-array and per-layer pulse distributions and the hottest cells, not only
      averages; the TT family's $A$ arrays and the resets dominate the write load, and the pulse budget per update clips the
      largest layers (Section~\ref{sec:main}).
\end{itemize}

\paragraph{Stage 2: the symmetry point and the reset as physical processes.}
\begin{itemize}
\item \textbf{Measurement.} The pulse count, latency and energy to bring a device to its symmetry point or to its read zero; the
      residual error, its spread and its drift over hours and $10^6$ cycles; the signed range left on each side.
\item \textbf{Why.} With a realistic pulse budget the reset is incomplete (on the protonic-inspired device, 62 alternating
      pulses leave \RVresetResidPct\,\% of a displacement), so the reset is a device specification, not a free operation.
\end{itemize}

\paragraph{Stage 3: small arrays with hardware in the loop.}
\begin{itemize}
\item \textbf{Hardware.} Two $64\times64$ arrays per layer ($A$ and $C$); forward and backward passes on the host from the read
      weights; updates on chip as stochastic pulse coincidences; resets or chopping as the algorithm requires.
\item \textbf{Prediction first.} Before training, fit the Stage~1 model on several independent chips, run this simulator with
      the fit, and state a \emph{prediction interval} that includes device-population variation, fit uncertainty, training
      randomness and measurement noise (not a confidence interval of the simulated mean, which shrinks with more seeds). The
      hardware should land inside it; independently, within 2\,pp of FP-SGD shows the device works.
\end{itemize}

\paragraph{Stage 4: full arrays and in-memory reads.} Move the forward and backward reads on chip and measure the system energy
per weight per row, with the read precision that gives the accuracy, against the conditional projection of
Section~\ref{sec:energy}; convert retention requirements into time with the full step duration.

\paragraph{The 1960 memistor.} Appendix~\ref{app:bench} gives a proposed, unbuilt bench replication protocol with hobby parts and
safety notes. Before its fits are used as metrology it must measure the actual read current, keep per-pulse charge, separate
read from write noise and reject incomplete pulses.

\section{Conclusion}
In this simulator, with ideal reads and every update made of device pulses, two electrochemical specifications (the AIHWKIT
Li-ECRAM fit and a protonic-inspired scenario) train three small networks to within the pre-set margins of SGD with exact updates
when a chopped Tiki-Taka algorithm moves the gradient; plain SGD at the same learning rate does not on any published device. Against
a tuned digital Adam none of them does, and on the language model the best is \RVLmAdamBestGap\,\% worse in bits per byte. No
growth of the gap with width was detected over a factor of four, without establishing equivalence. The energy advantage is a
conditional projection (\RVenSysMin--\RVenChipMax$\times$ per weight-row at borrowed inference-chip read efficiencies), reached
only where the quality is. After an independent review we repaired the simulator and report how much each repair moved the
original results. The remaining questions are experimental: the reset, the reference and the read path in real arrays.
